\newpage
\begin{center}
	\textbf{\large 1. Задача}
\end{center}
\refstepcounter{chapter}
\addcontentsline{toc}{chapter}{1. Задача}

\section{Краткое описание}

Целью курсовой работы является разработка приложения для локализации QR-кода на цифровом изображении большого разрешения с использованием библиотеки OpenCV и языка программирования C++.
Под локализацией понимается определение области изображения, занимаемой QR-кодом, в виде четырёх точек (вершин) многоугольника, аппроксимирующего границы QR-кода (без учёта \textit{quiet zone}).

Практическая мотивация задачи связана с инвентаризацией оборудования в дата-центре: QR-коды нанесены на инвентарные метки, закреплённые на серверах и сетевых устройствах, расположенных в стойках (rack).
При фотографировании оборудования общий кадр, как правило, содержит большое количество деталей стойки, при этом QR-код занимает небольшую часть изображения.
В результате прямое применение стандартных методов детектирования QR-кодов на изображении целиком может давать неудовлетворительные результаты.
Поэтому возникает необходимость предварительной локализации области QR-кода, позволяющей выделить компактный фрагмент изображения, на котором последующее распознавание становится более надёжным (см. рис.~\ref{nmp1} и \ref{nmp2}).

\begin{figure}[!h]
	\begin{center}
		\includegraphics[width=250pt]{im01.png}
		\caption{Функция \texttt{cv::QRCodeDetector} не справилась с детектированием QR-кода при подаче изображения целиком}
		\label{nmp1}
	\end{center}
\end{figure}

\begin{figure}[!h]
	\begin{center}
		\includegraphics[width=100pt]{im02.png}
		\caption{После локализации \texttt{cv::QRCodeDetector} успешно справилась с детекцией и декодированием}
		\label{nmp2}
	\end{center}
\end{figure}

\setcounter{page}{3}

\section{Подробное описание}
Входными данными является путь к изображению, содержащему один QR-код на однородном (преимущественно белом) фоне.
Практически это соответствует инвентарным меткам, применяемым для учёта оборудования: QR-код печатается на светлой этикетке вместе с инвентарным номером, и при фотографировании в условиях дата-центра занимает небольшую часть кадра на фоне элементов стойки и корпуса устройства.

Изображения могут включать поворот QR-кода на произвольный угол относительно осей изображения, а также небольшие перспективные искажения, обусловленные нестрогой перпендикулярностью камеры к плоскости метки.
При этом компенсация сложных факторов съёмки (неравномерное освещение, сильные блики, тени, выраженный шум матрицы, цветовые искажения) не входила в постановку задачи: предполагалось, что фон метки однороден, а качество исходных изображений достаточно для выделения геометрических признаков QR-кода.

В качестве результата алгоритм должен вернуть четыре точки $\{(x_i, y_i)\}_{i=1}^{4}$,
соответствующие вершинам многоугольника, аппроксимирующего границы QR-кода (без включения \textit{quiet zone}).
Порядок точек фиксирован: \texttt{TL} (верхний левый), \texttt{TR} (верхний правый), \texttt{BR} (нижний правый), \texttt{BL} (нижний левый).
Порядок обхода вершин многоугольника показан на рис.~\ref{nmp3}.

Для количественной оценки качества была сформирована эталонная разметка (Ground Truth).
Разметка создавалась вручную с использованием инструмента LabelMe: для каждого изображения выделялась область QR-кода в виде четырёхугольника, соответствующего реальным границам кода.

Критерии качества были выбраны следующими:

\begin{itemize}
	\item Корректное определение факта наличия QR-кода на изображении;
	\item Высокая точность границ (оценивается метрикой IoU при сравнении с Ground Truth, а также ошибкой по углам);
	\item Устойчивость к масштабу (QR-код может быть маленьким относительно изображения большого разрешения);
	\item Приемлемая вычислительная сложность при обработке изображений большого размера.
\end{itemize}

\begin{figure}[!h]
	\begin{center}
		\includegraphics[width=200pt]{im03.png}
		\caption{Порядок обхода вершин многоугольника: \texttt{TL}, \texttt{TR}, \texttt{BR}, \texttt{BL}}
		\label{nmp3}
	\end{center}
\end{figure}

\newpage

Формат JSON файла, который служит результатом работы приложения, следующий:
\begin{verbatim}
{
  "shapes": [
    {
      "label": "qr",
      "points": [
        [x_tl, y_tl],
        [x_tr, y_tr],
        [x_br, y_br],
        [x_bl, y_bl]
      ],
      "description": "<описание_изображение_для_отладки>",
      "shape_type":  "polygon"
    }
  ],
  "imagePath":  "<имя_файла_изображения>",
  "imageHeight": <высота_изображения_в_пикселях>,
  "imageWidth":  <ширина_изображения_в_пикселях>
}
\end{verbatim}

К технической реализации приложения предъявляются следующие требования:
\begin{itemize}
	\item Язык программирования C++, версия не ниже 17;
	\item Сборка с использованием CMake, версия не ниже 3.15;
	\item Реализация с использованием библиотек \texttt{OpenCV} и \texttt{nlohmann\_json};
	\item Кроссплатформенная сборка (под разные операционные системы);
	\item Реализация взаимодействия с пользователем через консольный интерфейс.
\end{itemize}

\begin{figure}[!h]
	\begin{center}
		\includegraphics[width=250pt]{im04.jpg}
		\caption{Изображение, подающееся на вход программы}
		\label{nmp4}
	\end{center}
\end{figure}

\begin{figure}[!h]
	\begin{center}
		\includegraphics[width=250pt]{im05.png}
		\caption{Результат работы программы (визуализация)}
		\label{nmp5}
	\end{center}
\end{figure}
